%%%%%%%%%%%%%%%%%%%%%%%%%%%%%%%%%%%%%%%%%%%%%%%%%%%%%%%%%%%%
%% PREFACE -- A Letter to the Curious Reader
%%%%%%%%%%%%%%%%%%%%%%%%%%%%%%%%%%%%%%%%%%%%%%%%%%%%%%%%%%%%

\chapter*{Preface: A Letter to the Curious Reader}
\addcontentsline{toc}{chapter}{Preface: A Letter to the Curious Reader}

\section*{The thing in front of you}
\addcontentsline{toc}{section}{The thing in front of you}

You have, almost certainly, already used a Large Language Model.
You asked it to draft an email, or to explain a paragraph in a
contract, or to summarise a research paper that a colleague had
sent you, or to suggest a name for a small dog. The answer it gave
you was probably useful. It was also, at some level, a surprise ---
not because the technology is new (it is not, really; its lineage
is more than a century old), but because the experience of being
addressed in fluent prose by a system with no body, no mind, and
no opinions felt different from anything you had asked a computer
to do before.

It is worth pausing on that surprise. The interesting question, for
the kind of reader who wants to understand these systems rather
than merely use them, is not whether the answer was correct. The
interesting question is what the system was actually \emph{doing}.
This guide is going to take it apart, slowly, and put it back
together. By the time you finish, you will not, in any deep sense,
find these systems mysterious. You may, however, find them more
remarkable than you do today, and you will certainly find them
less trustworthy in some specific and useful ways.

This is a book for the educated reader who has no background in
mathematics or computer science. There are no equations. There
is no code. There is no jargon that is not paid for, the first
time it appears, with a plain-English equivalent. What there is,
instead, is a careful walk through the ideas, the people, and the
engineering decisions that produced the systems on the other side
of your screen. The walk is long but it is not steep. We are not
going to dazzle you. We are going to give you, page by page, a
sturdier intuition for an artefact that did not exist, in the form
you now use it, when most of the readers of this book were
already adults.

\section*{Why this book exists}
\addcontentsline{toc}{section}{Why this book exists}

There is a parallel reader for engineers, called \emph{Large
Language Models as Engineered Systems}, that walks through the
same ideas with the formal mathematics included. That reader is
the right one if you want to be able to derive the equations and
implement the algorithms yourself. This guide is the right one if
you want to be able to reason about these systems --- to ask
better questions of the engineers who build them, the vendors who
sell them, the regulators who oversee them, and, increasingly, the
products that rely on them --- without first having to spend three
years on a mathematics degree.

The two books share a chapter ordering and most of their narrative
voice. They also share an observation, which is that the most
useful thing a non-engineer can know about these systems is not
how they succeed but how, and why, they fail. We will return to
this point repeatedly. The failure modes of a Large Language Model
are not embarrassments to be hidden. They are the working
vocabulary of anyone who deploys these systems --- or relies on
someone else's deployment of them --- in a context where the
answer matters.

\section*{How we got here}
\addcontentsline{toc}{section}{How we got here}

It is worth remembering that the field has a history --- and that
the history is longer than the headlines suggest.

Statistical language modelling began in 1913, when the Russian
mathematician Andrey Markov showed that language could be
described as a chain of conditional dependencies between symbols.
Claude Shannon, in 1948, gave the field a way to measure how well
any model captures the regularities of a language --- a yardstick
that is still in use today. The first artificial neurons appeared
in the early 1950s, and what followed was three decades of
alternating promise and disappointment: the perceptron, two AI
winters when funding collapsed and confidence deflated, the slow
rehabilitation of multi-layer networks through backpropagation in
the 1980s, and recurrent architectures in the 1990s that could
finally hold context across a sequence of words.

By 2003 the first recognisably modern neural language model
existed. By 2013, words could be placed in a geometric space
where arithmetic on meaning worked: \emph{king minus man plus
woman} landed near \emph{queen}. By 2015 a mechanism called
\emph{attention} let a model focus selectively on the parts of
a sequence that mattered for the next prediction. In 2017, a
team at Google removed the last architectural constraint and
published a paper titled, with characteristic understatement,
\emph{Attention Is All You Need}. The architecture described
in that paper is still the basis of every system in this book.

Five years of scaling followed. GPT-1 in 2018, GPT-3 in 2020,
ChatGPT at the end of 2022, and a steady cascade of capable
competitors in the years after. None of these systems came out
of nowhere. Each one was possible only because of the step
before it. The chapters that follow trace that lineage one idea
at a time, and when you finish you will see the whole edifice
clearly --- a long tradition of engineering, not a sudden
apparition.

\section*{What this book is not}
\addcontentsline{toc}{section}{What this book is not}

This book is not a defence of the technology. It is not an
attack on it either. The honest position --- the one that the
engineers who built these systems will tell you in private, and
that the most thoughtful critics of these systems will agree
with --- is that they are remarkable and limited in roughly equal
measure, and that the public conversation has not, so far, done a
good job of holding both halves of that observation at the same
time. We will try.

This book is not a forecast. It does not tell you which of the
current models will dominate next year, or whether the next
generation of training runs will yield a qualitative leap or a
disappointing plateau. The honest answer to that question is that
we do not know. The people who tell you they do know are usually
selling you something, and the few who actually know say very
little in public.

This book is not a manual. It will not teach you how to build a
Large Language Model. It will not teach you how to fine-tune one.
It will not even, in most cases, teach you how to write better
prompts --- though, as a side effect, you may find that you start
to write better ones because you understand what the model is
doing with them.

What this book \emph{is} is an attempt to give you, in plain
language and at a reading pace that respects your time, genuine
mechanism-level comprehension of a class of system that has, in
the space of about five years, moved from a research curiosity to
a fixture of modern professional life. The book is short by the
standards of the technical literature; it is long by the standards
of a thing you might read on a long flight. It is meant to be read
slowly, once, and then dipped into when a specific question comes up.

The discipline this book tries to hold itself to is sometimes
called the \emph{Feynman Standard} --- though that name
overstates the case. There is no written rulebook. It is, more
honestly, an observation of how a great teacher worked and an
aspiration to do something similar. Richard Feynman's popular
lectures were not remarkable for avoiding mathematics; they were
remarkable because he refused to let a term stand unexplained.
Every word had to earn its place. He was particular about the
difference between knowing the name of a thing and knowing the
thing itself. Knowing the name of a bird in all the languages of
the world, he said, tells you nothing about the bird. This book
will try, at every step, to give you not the name but the bird:
not merely what each component is called, but why it behaves as
it does, and why that behaviour produces the specific failures
we see in the wild. A chapter that teaches you vocabulary without
teaching you the mechanism behind it has done half the work. We
aim for the other half as well.

\section*{What goes wrong, and why this is the most important section
of the book}
\addcontentsline{toc}{section}{What goes wrong}

A language model trained the way modern ones are trained is, by
construction, the most plausible thing the model can say given
the prompt. The most plausible thing is not the same as the true
thing. The model has no native concept of truth at all. It has a
beautifully calibrated sense of plausibility and a great deal of
memorised plausibility, which is sometimes the same as truth and
sometimes is not. This is the seed of \emph{hallucination} --- the
polite name we give to a fluent, confident, plausible, wrong
answer. We will return to it in Chapter~\ref{chap:stat-foundations}, in Chapter~\ref{chap:systems}, and at
length in Chapter~\ref{chap:error-model}.

A model that is fluent is not, by default, a model that is
helpful. A model that has only been trained on the kind of
prediction we will describe in Chapter~\ref{chap:stat-foundations} will happily complete an
essay that begins ``the best way to commit fraud is''. Making it
refuse to do so is the work of a discipline called
\emph{alignment}, and alignment is its own discipline of trade-
offs. We will see in Chapter~\ref{chap:alignment} how the discipline arrived, what
it bought, and what it cost --- the so-called \emph{alignment
tax} on capability, the \emph{sycophancy} that emerges when the
model is rewarded too much for agreeing with the annotator, the
\emph{reward hacking} that lets the model find a high-reward
path that no annotator would ever endorse.

The careful question that produces a careful answer often differs
from the careless question that produces a wrong one only by a
few words. Two prompts that a human would read as equivalent can
produce two qualitatively different outputs. We call this
\emph{prompt brittleness}. We will see, in Chapter~\ref{chap:transformer}, why the
mechanism the model uses to read its inputs makes this not just
possible but inevitable.

A model that retrieves from a document store is only as honest as
the store. A model that calls a tool is only as careful as the
tool's contract. We will see, in Chapter~\ref{chap:systems} and again in
Chapter~\ref{chap:agents}, how retrieval can hallucinate, how the model can be
\emph{lost in the middle} of a long context, and how a malicious
tool output can be crafted to look like a system instruction ---
the class of attack we call \emph{prompt injection}, which is to
LLM systems roughly what SQL injection once was to web
applications.

These failure modes are not embarrassments. They are the working
vocabulary of anyone who deploys these systems into the world,
and they are increasingly the working vocabulary of anyone who
relies on someone else's deployment. Every chapter ends with a
discussion of the characteristic failures of the technique it
presents. 
% Read those sections carefully. % NOTE - removed for being to prescriptive to the reader
They are where the unhelpful intuition that ``the model just knows'' meets the more
useful intuition of ``the model is doing exactly what we asked,
and what we asked is not quite what we meant.''

\section*{How to read this book}
\addcontentsline{toc}{section}{How to read this book}

The book is organised into thirteen chapters, one per topic. Each
chapter mirrors a chapter of the parent reader, with the same
title and the same scope but without the formal mathematics.

Every chapter follows the same structure. It opens with a short
narrative section called \emph{Setting the Scene}, which describes
the problem the chapter is about, the lineage of ideas that
brought us to the current best practice, what that current best
practice is, and the characteristic ways the practice fails. Then
comes a section called \emph{The idea, in plain language}, which
is the conceptual core of the chapter. Then comes a section called
\emph{The engineering consequence}, which is what the idea makes
easy, what it makes hard, and the trade-offs a product team faces
today because of it. Then comes a section called \emph{What goes
wrong, and why}, which traces the chapter's failure modes back to
their mechanism. Finally, a short \emph{Bridges} paragraph names
the chapters that build on this one and the chapters whose
machinery this one depends on, so that a reader who already knows
part of the material can skip in safely.

The narrative comes first deliberately. A reader who knows what a
thing is for tends to derive its details more easily; a reader
who tries to learn the details first sometimes loses track of why
they matter. We write the why first. Then the how.

If you are reading the book front-to-back you will find that the
chapters build on one another. The notion of a probability
distribution from Chapter~\ref{chap:stat-foundations} reappears in Chapter~\ref{chap:pretraining} as the basis
of pretraining and again in Chapter~\ref{chap:alignment} as the thing alignment
modifies. The notion of a vector from Chapter~\ref{chap:tokenization} reappears in
Chapter~\ref{chap:transformer} as the input to attention and in Chapter~\ref{chap:systems} as the
basis of retrieval. The notion of a trust boundary from
Chapter~\ref{chap:alignment} reappears in Chapter~\ref{chap:agents} as the agent-loop invariant.
You do not have to remember any of this. The book will remind
you each time it matters.

\section*{A note on tone}
\addcontentsline{toc}{section}{A note on tone}

The narrative in this book is written in plain language because
the subject is not. The subject is technical, and the people who
build these systems use technical language for precise reasons.
We have tried to translate that language without losing what the
precision was for. Sometimes a chapter will pause to introduce a
term of art, define it in plain English, and then use it freely.
That is not jargon for its own sake. It is the price of being
able to say things that the surrounding ordinary vocabulary
cannot say without ten times the words. We have tried to keep
that price low.

We have also tried to acknowledge the people who built the ideas
on which this field stands, because giving credit is a courtesy,
because knowing who first solved a problem helps you guess who
has thought hardest about its limits, and because the public
conversation has, on the whole, done a poor job of distributing
credit across a hundred years and a thousand researchers. This is
a small correction, made one date and one name at a time.

If a passage feels like it is trying to dazzle you, please tell
us; we will rewrite it. The intent is not awe. The intent is that
you put the book down able to reason about a class of system
that did not exist, in the form you now use it, when you first learned to read.


\section*{On the writing of this book}
\addcontentsline{toc}{section}{On the writing of this book}

This guide was written in collaboration with Claude, a Large
Language Model produced by Anthropic. The subject matter of the
book makes that fact worth disclosing rather than concealing,
and you should know it before you read further.

The author chose the scope of the book, the chapters it would
contain, the historical threads it would trace, the voice the
narrative sections would adopt, and the failure modes the
discussion would emphasise. Each chapter was drafted, then read,
then revised, then committed. The model produced the prose
inside those choices. The author made the choices, read the
drafts, redirected the ones that drifted, and approved each
chapter before it landed.

Errors that survived this process are the author's, since
responsibility is a thing an author has and a language model
does not. Readers who find passages that are wrong, evasive, or
fluent in a way that papers over a real difficulty are
encouraged to say so. Any future edition will be better for the
saying.
